\chapter{理论力学回顾：支座反力}

\section{平衡方程速记}

\begin{notecard}
求支座反力时，先把结构作为整体隔离。能用一个力矩方程直接消去未知反力时，优先对该支座或铰点取矩；再用
\[
  \sum F_x=0,\qquad \sum F_y=0,\qquad \sum M=0
\]
补齐其余反力。分布荷载先等效为合力：均布荷载 \(q\) 作用长度 \(L\) 时合力为 \(qL\)，作用在荷载段中点；三角形荷载最大值为 \(q_0\)、底边为 \(l\) 时合力为 \(\frac12q_0l\)，作用在距最大端 \(l/3\) 处。若算得某反力为负，表示实际方向与预先假定方向相反。
\end{notecard}

\section{第 2 次作业：求支座反力}

\begin{problemcard}{2-1(a) 三角形分布荷载悬臂梁}
求图示结构的支座反力。梁 \(AB\) 长为 \(l\)，右端 \(B\) 固定，三角形分布荷载从 \(A\) 端为零线性增大到 \(B\) 端的 \(q_0\)。

\begin{figure}[H]
\centering
\begin{tikzpicture}[scale=0.95,>=Stealth]
  \draw[member] (0,0)--(4,0);
  \draw[ground] (4,-0.55)--(4,0.55);
  \foreach \y in {-0.42,-0.24,-0.06,0.12,0.30,0.48}
    \draw[ground] (4,\y)--(4.18,\y+0.10);
  \foreach \x/\h in {0.55/0.20,1.05/0.34,1.55/0.48,2.05/0.62,2.55/0.76,3.05/0.90,3.55/1.04}
    \draw[Brand,-{Stealth[length=1.8mm]}] (\x,\h)--(\x,0.08);
  \draw[Brand] (0.45,0.20)--(3.75,1.12);
  \draw[<->] (0,-0.42)--node[below] {\(l\)} (4,-0.42);
  \draw[<->] (4.25,0)--node[right] {\(q_0\)} (4.25,1.12);
  \node[below] at (0,0) {\(A\)};
  \node[below left] at (4,0) {\(B\)};
  \figtitle{(a) 固定端承受三角形分布荷载}
\end{tikzpicture}
\end{figure}
\end{problemcard}

\begin{solutioncard}
\textbf{答案：}\answer{\(F_{By}=\dfrac12q_0l\) 向上，\(M_B=\dfrac16q_0l^2\) 顺时针。}

\textbf{解析：}三角形荷载等效为一个向下合力
\[
  W=\frac12q_0l .
\]
合力作用点在距荷载强度最大端 \(B\) 为 \(l/3\) 的位置。由整体竖向力平衡，
\[
  \sum F_y=0:\qquad F_{By}-W=0
  \quad\Rightarrow\quad
  F_{By}=\frac12q_0l .
\]
对 \(B\) 点取矩，三角形荷载对 \(B\) 的力臂为 \(l/3\)，因此固定端反力矩大小为
\[
  M_B=W\cdot\frac{l}{3}
  =\frac12q_0l\cdot\frac{l}{3}
  =\frac16q_0l^2 .
\]
该反力矩方向与外荷载对 \(B\) 的转动趋势相反，按题解记为顺时针。
\end{solutioncard}

\begin{problemcard}{2-1(b) 集中力与局部均布荷载}
悬臂梁 \(AB\) 左端 \(A\) 固定，全长 \(3\,\mathrm{m}\)。右端 \(B\) 作用 \(30\,\mathrm{kN}\) 向下集中力，末端 \(CB\) 长 \(1\,\mathrm{m}\)，承受 \(15\,\mathrm{kN/m}\) 均布荷载。求固定端反力。

\begin{figure}[H]
\centering
\begin{tikzpicture}[scale=0.95,>=Stealth]
  \draw[member] (0,0)--(4.5,0);
  \draw[ground] (0,-0.50)--(0,0.50);
  \foreach \y in {-0.40,-0.22,-0.04,0.14,0.32}
    \draw[ground] (0,\y)--(-0.18,\y-0.10);
  \foreach \x in {3.15,3.45,3.75,4.05,4.35}
    \draw[Brand,-{Stealth[length=1.8mm]}] (\x,0.62)--(\x,0.08);
  \draw[Brand] (3.05,0.62)--(4.45,0.62);
  \draw[Brand,-{Stealth[length=2mm]}] (4.5,1.10)--(4.5,0.08);
  \node[above] at (3.75,0.70) {\(15\,\mathrm{kN/m}\)};
  \node[right] at (4.5,1.00) {\(30\,\mathrm{kN}\)};
  \draw[<->] (0,-0.42)--node[below] {\(3\,\mathrm{m}\)} (4.5,-0.42);
  \draw[<->] (3.0,-0.78)--node[below] {\(1\,\mathrm{m}\)} (4.5,-0.78);
  \node[below] at (0,0) {\(A\)};
  \node[below] at (3.0,0) {\(C\)};
  \node[below] at (4.5,0) {\(B\)};
  \figtitle{(b) 固定端梁}
\end{tikzpicture}
\end{figure}
\end{problemcard}

\begin{solutioncard}
\textbf{答案：}\answer{\(F_{Ay}=45\,\mathrm{kN}\) 向上，\(M_A=127.5\,\mathrm{kN\cdot m}\) 逆时针。}

\textbf{解析：}末端均布荷载的合力为
\[
  15\times 1=15\,\mathrm{kN},
\]
作用在 \(CB\) 中点，距 \(A\) 为 \(2.5\,\mathrm{m}\)。竖向平衡得
\[
  F_{Ay}-15\times1-30=0
  \quad\Rightarrow\quad
  F_{Ay}=45\,\mathrm{kN}.
\]
对 \(A\) 点取矩，外荷载使梁产生顺时针转动，固定端反力矩取相反方向：
\[
  M_A=15\times1\times2.5+30\times3
  =37.5+90
  =127.5\,\mathrm{kN\cdot m}.
\]
\end{solutioncard}

\begin{problemcard}{2-1(c) 带外伸段的简支梁}
梁上 \(A\) 端作用 \(30\,\mathrm{kN\cdot m}\) 力偶，\(C\)、\(B\) 为竖向支座；\(AC=1\,\mathrm{m}\)，\(AB=4\,\mathrm{m}\)，集中力 \(40\,\mathrm{kN}\) 作用于 \(D\)，且 \(DB=1.5\,\mathrm{m}\)。求 \(C\)、\(B\) 处支反力。

\begin{figure}[H]
\centering
\begin{tikzpicture}[scale=0.95,>=Stealth]
  \draw[member] (0,0)--(5.0,0);
  \roller{1.25}{0}
  \roller{5.0}{0}
  \draw[Brand,-{Stealth[length=2mm]}] (3.15,1.0)--(3.15,0.08);
  \node[right] at (3.15,0.86) {\(40\,\mathrm{kN}\)};
  \draw[Brand,-{Stealth[length=2mm]}] (0.20,0.55) arc[start angle=65,end angle=300,radius=0.36];
  \node[above left] at (0.15,0.55) {\(30\,\mathrm{kN\cdot m}\)};
  \draw[<->] (0,-0.55)--node[below] {\(1\,\mathrm{m}\)} (1.25,-0.55);
  \draw[<->] (0,-0.90)--node[below] {\(4\,\mathrm{m}\)} (5,-0.90);
  \draw[<->] (3.15,-0.55)--node[below] {\(1.5\,\mathrm{m}\)} (5,-0.55);
  \node[below] at (0,0) {\(A\)};
  \node[below] at (1.25,0) {\(C\)};
  \node[below] at (3.15,0) {\(D\)};
  \node[below] at (5,0) {\(B\)};
  \figtitle{(c) 对 \(C\) 点取矩可直接求 \(B\) 反力}
\end{tikzpicture}
\end{figure}
\end{problemcard}

\begin{solutioncard}
\textbf{答案：}\answer{\(F_B=10\,\mathrm{kN}\) 向上，\(F_C=30\,\mathrm{kN}\) 向上。}

\textbf{解析：}取整体隔离体。对 \(C\) 点取矩时，\(F_C\) 的力矩为零，\(B\) 到 \(C\) 的距离为 \(3\,\mathrm{m}\)，\(40\,\mathrm{kN}\) 到 \(C\) 的距离为 \(1.5\,\mathrm{m}\)。按题解的符号约定列式：
\[
  F_B\times3+30-40\times1.5=0 .
\]
故
\[
  F_B=\frac{40\times1.5-30}{3}=10\,\mathrm{kN}.
\]
再由竖向力平衡
\[
  F_C+F_B-40=0
  \quad\Rightarrow\quad
  F_C=40-10=30\,\mathrm{kN}.
\]
\end{solutioncard}

\begin{problemcard}{2-1(d) 均布荷载与集中力偶}
简支梁 \(AB\) 全长 \(10\,\mathrm{m}\)，全梁受 \(0.2\,\mathrm{kN/m}\) 均布荷载；\(C\) 点距 \(B\) 为 \(2\,\mathrm{m}\)，在 \(C\) 点作用 \(4\,\mathrm{kN\cdot m}\) 力偶。求两端支反力。

\begin{figure}[H]
\centering
\begin{tikzpicture}[scale=0.88,>=Stealth]
  \draw[member] (0,0)--(6,0);
  \roller{0}{0}
  \roller{6}{0}
  \foreach \x in {0.25,0.65,...,5.85}
    \draw[Brand,-{Stealth[length=1.7mm]}] (\x,0.62)--(\x,0.08);
  \draw[Brand] (0.15,0.62)--(5.95,0.62);
  \draw[Brand,-{Stealth[length=2mm]}] (4.80,0.85) arc[start angle=120,end angle=-160,radius=0.36];
  \node[above] at (2.25,0.80) {\(0.2\,\mathrm{kN/m}\)};
  \node[above right] at (4.80,0.98) {\(4\,\mathrm{kN\cdot m}\)};
  \draw[<->] (0,-0.50)--node[below] {\(10\,\mathrm{m}\)} (6,-0.50);
  \draw[<->] (4.8,-0.86)--node[below] {\(2\,\mathrm{m}\)} (6,-0.86);
  \node[below] at (0,0) {\(A\)};
  \node[below] at (4.8,0) {\(C\)};
  \node[below] at (6,0) {\(B\)};
  \figtitle{(d) 全跨均布荷载}
\end{tikzpicture}
\end{figure}
\end{problemcard}

\begin{solutioncard}
\textbf{答案：}\answer{\(F_B=1.4\,\mathrm{kN}\) 向上，\(F_A=0.6\,\mathrm{kN}\) 向上。}

\textbf{解析：}全梁均布荷载合力为
\[
  W=0.2\times10=2\,\mathrm{kN},
\]
作用在跨中，距 \(A\) 为 \(5\,\mathrm{m}\)。对 \(A\) 点取矩：
\[
  F_B\times10-4-0.2\times10\times5=0 .
\]
于是
\[
  F_B=\frac{4+10}{10}=1.4\,\mathrm{kN}.
\]
竖向力平衡给出
\[
  F_A+F_B-2=0
  \quad\Rightarrow\quad
  F_A=2-1.4=0.6\,\mathrm{kN}.
\]
该题的关键是把集中力偶直接作为力矩项加入方程，不需要给它再乘距离。
\end{solutioncard}

\begin{problemcard}{2-1(e) 两端力偶与集中力}
简支梁 \(AB\) 长 \(3\,\mathrm{m}\)，\(A\) 端作用 \(6\,\mathrm{kN\cdot m}\) 力偶，\(B\) 端作用 \(20\,\mathrm{kN\cdot m}\) 力偶；\(C\) 点距 \(B\) 为 \(1\,\mathrm{m}\)，有 \(20\,\mathrm{kN}\) 集中力向下。求支反力。

\begin{figure}[H]
\centering
\begin{tikzpicture}[scale=0.95,>=Stealth]
  \draw[member] (0,0)--(4.2,0);
  \roller{0}{0}
  \roller{4.2}{0}
  \draw[Brand,-{Stealth[length=2mm]}] (0.25,0.62) arc[start angle=60,end angle=300,radius=0.35];
  \draw[Brand,-{Stealth[length=2mm]}] (4.00,0.62) arc[start angle=120,end angle=-160,radius=0.35];
  \draw[Brand,-{Stealth[length=2mm]}] (2.8,0.95)--(2.8,0.08);
  \node[above left] at (0.30,0.72) {\(6\,\mathrm{kN\cdot m}\)};
  \node[above right] at (3.95,0.98) {\(20\,\mathrm{kN\cdot m}\)};
  \node[left] at (2.75,1.05) {\(20\,\mathrm{kN}\)};
  \draw[<->] (0,-0.52)--node[below] {\(3\,\mathrm{m}\)} (4.2,-0.52);
  \draw[<->] (2.8,-0.88)--node[below] {\(1\,\mathrm{m}\)} (4.2,-0.88);
  \node[below] at (0,0) {\(A\)};
  \node[below] at (2.8,0) {\(C\)};
  \node[below] at (4.2,0) {\(B\)};
  \figtitle{(e) 集中力偶不乘距离}
\end{tikzpicture}
\end{figure}
\end{problemcard}

\begin{solutioncard}
\textbf{答案：}\answer{\(F_B=22\,\mathrm{kN}\) 向上；若向上为正，\(F_A=-2\,\mathrm{kN}\)，即 \(A\) 处实际向下，大小为 \(2\,\mathrm{kN}\)。}

\textbf{解析：}集中力 \(20\,\mathrm{kN}\) 作用点距 \(A\) 为 \(2\,\mathrm{m}\)。对 \(A\) 点取矩，按题解中力偶方向写为
\[
  F_B\times3-6-20\times2-20=0 .
\]
解得
\[
  F_B=\frac{6+40+20}{3}=22\,\mathrm{kN}.
\]
由竖向平衡
\[
  F_A+F_B-20=0
  \quad\Rightarrow\quad
  F_A=20-22=-2\,\mathrm{kN}.
\]
因此若以向上为正，\(F_A=-2\,\mathrm{kN}\)；题解图按反力实际方向标注时，可记为 \(A\) 处 \(2\,\mathrm{kN}\) 向下。若只列反力大小，应同时注明方向，避免把符号含义混淆。
\end{solutioncard}

\begin{problemcard}{2-1(f) 伸臂梁端部均布荷载}
梁 \(AB\) 总长 \(5a\)，在 \(A\) 与 \(C\) 处支承，\(CB=a\)。伸臂段 \(CB\) 上作用强度为 \(q\) 的均布荷载。求 \(A\)、\(C\) 两处支反力。

\begin{figure}[H]
\centering
\begin{tikzpicture}[scale=0.88,>=Stealth]
  \draw[member] (0,0)--(6,0);
  \roller{0}{0}
  \roller{4.8}{0}
  \foreach \x in {5.0,5.25,5.50,5.75}
    \draw[Brand,-{Stealth[length=1.7mm]}] (\x,0.62)--(\x,0.08);
  \draw[Brand] (4.85,0.62)--(6.0,0.62);
  \node[above] at (5.4,0.72) {\(q\)};
  \draw[<->] (0,-0.52)--node[below] {\(5a\)} (6,-0.52);
  \draw[<->] (4.8,-0.88)--node[below] {\(a\)} (6,-0.88);
  \node[below] at (0,0) {\(A\)};
  \node[below] at (4.8,0) {\(C\)};
  \node[below] at (6,0) {\(B\)};
  \figtitle{(f) 右端伸臂段受均布荷载}
\end{tikzpicture}
\end{figure}
\end{problemcard}

\begin{solutioncard}
\textbf{答案：}\answer{\(F_C=\dfrac98qa\) 向上；若向上为正，\(F_A=-\dfrac18qa\)，即 \(A\) 处实际向下，大小为 \(\dfrac18qa\)。}

\textbf{解析：}伸臂段上的均布荷载合力为 \(qa\)，作用点在 \(CB\) 中点。由题图尺寸，\(AC=4a\)，合力作用点距 \(A\) 为
\[
  4a+\frac{a}{2}=4.5a .
\]
对 \(A\) 点取矩：
\[
  F_C\cdot4a-qa\cdot4.5a=0
  \quad\Rightarrow\quad
  F_C=\frac{4.5}{4}qa=\frac98qa .
\]
再由竖向力平衡
\[
  F_A+F_C-qa=0
  \quad\Rightarrow\quad
  F_A=qa-\frac98qa=-\frac18qa .
\]
负号说明 \(A\) 处反力实际方向向下，这也是伸臂梁端部荷载常见的“压起”现象。
\end{solutioncard}

\begin{problemcard}{2-1(g) 局部均布荷载悬臂梁}
悬臂梁 \(AB\) 左端 \(A\) 固定，全长 \(2.5\,\mathrm{m}\)，右端 \(CB\) 长 \(1\,\mathrm{m}\)，承受 \(5\,\mathrm{kN/m}\) 均布荷载。求固定端反力。

\begin{figure}[H]
\centering
\begin{tikzpicture}[scale=0.95,>=Stealth]
  \draw[member] (0,0)--(4.5,0);
  \draw[ground] (0,-0.50)--(0,0.50);
  \foreach \y in {-0.40,-0.22,-0.04,0.14,0.32}
    \draw[ground] (0,\y)--(-0.18,\y-0.10);
  \foreach \x in {2.75,3.10,3.45,3.80,4.15}
    \draw[Brand,-{Stealth[length=1.7mm]}] (\x,0.62)--(\x,0.08);
  \draw[Brand] (2.7,0.62)--(4.5,0.62);
  \node[above] at (3.55,0.72) {\(5\,\mathrm{kN/m}\)};
  \draw[<->] (0,-0.50)--node[below] {\(2.5\,\mathrm{m}\)} (4.5,-0.50);
  \draw[<->] (2.7,-0.86)--node[below] {\(1\,\mathrm{m}\)} (4.5,-0.86);
  \node[below] at (0,0) {\(A\)};
  \node[below] at (2.7,0) {\(C\)};
  \node[below] at (4.5,0) {\(B\)};
  \figtitle{(g) 合力距固定端 \(2\,\mathrm{m}\)}
\end{tikzpicture}
\end{figure}
\end{problemcard}

\begin{solutioncard}
\textbf{答案：}\answer{\(F_{Ay}=5\,\mathrm{kN}\) 向上，\(M_A=10\,\mathrm{kN\cdot m}\) 逆时针。}

\textbf{解析：}均布荷载合力为
\[
  W=5\times1=5\,\mathrm{kN}.
\]
荷载段长 \(1\,\mathrm{m}\)，其中点距 \(B\) 为 \(0.5\,\mathrm{m}\)，因此距固定端 \(A\) 为 \(2.5-0.5=2\,\mathrm{m}\)。竖向平衡得
\[
  F_{Ay}=5\,\mathrm{kN}.
\]
对 \(A\) 点取矩：
\[
  M_A=5\times2=10\,\mathrm{kN\cdot m}.
\]
反力矩方向与荷载产生的顺时针转动相反，故为逆时针。
\end{solutioncard}

\begin{problemcard}{2-1(h) 集中力与集中力偶悬臂梁}
悬臂梁 \(AB\) 左端固定，全长 \(2\,\mathrm{m}\)。\(C\) 点距 \(A\) 为 \(1\,\mathrm{m}\)，在 \(C\) 点作用 \(15\,\mathrm{kN}\) 向下集中力和 \(10\,\mathrm{kN\cdot m}\) 力偶。求固定端反力。

\begin{figure}[H]
\centering
\begin{tikzpicture}[scale=0.95,>=Stealth]
  \draw[member] (0,0)--(4.4,0);
  \draw[ground] (0,-0.50)--(0,0.50);
  \foreach \y in {-0.40,-0.22,-0.04,0.14,0.32}
    \draw[ground] (0,\y)--(-0.18,\y-0.10);
  \draw[Brand,-{Stealth[length=2mm]}] (2.2,1.05)--(2.2,0.08);
  \draw[Brand,-{Stealth[length=2mm]}] (2.45,0.80) arc[start angle=120,end angle=-160,radius=0.36];
  \node[above left] at (2.2,1.02) {\(15\,\mathrm{kN}\)};
  \node[above right] at (2.35,0.80) {\(10\,\mathrm{kN\cdot m}\)};
  \draw[<->] (0,-0.50)--node[below] {\(1\,\mathrm{m}\)} (2.2,-0.50);
  \draw[<->] (0,-0.86)--node[below] {\(2\,\mathrm{m}\)} (4.4,-0.86);
  \node[below] at (0,0) {\(A\)};
  \node[below] at (2.2,0) {\(C\)};
  \node[below] at (4.4,0) {\(B\)};
  \figtitle{(h) 集中力与集中力偶叠加}
\end{tikzpicture}
\end{figure}
\end{problemcard}

\begin{solutioncard}
\textbf{答案：}\answer{\(F_{Ay}=15\,\mathrm{kN}\) 向上，\(M_A=25\,\mathrm{kN\cdot m}\) 逆时针。}

\textbf{解析：}竖向外力只有 \(15\,\mathrm{kN}\) 向下，因此
\[
  F_{Ay}=15\,\mathrm{kN}.
\]
对 \(A\) 点取矩时，集中力贡献 \(15\times1=15\,\mathrm{kN\cdot m}\)，题给集中力偶贡献 \(10\,\mathrm{kN\cdot m}\)。所以固定端反力矩大小为
\[
  M_A=10+15\times1=25\,\mathrm{kN\cdot m}.
\]
方向为抵抗外荷载转动的逆时针方向。
\end{solutioncard}

\begin{problemcard}{2-1(i) 对称均布荷载简支梁}
简支梁 \(AB\) 全长 \(5a\)。中部 \(CD\) 长 \(3a\)，承受强度为 \(q\) 的均布荷载；\(AC=a\)，\(DB=a\)。求 \(A\)、\(B\) 两处支反力。

\begin{figure}[H]
\centering
\begin{tikzpicture}[scale=0.85,>=Stealth]
  \draw[member] (0,0)--(6,0);
  \roller{0}{0}
  \roller{6}{0}
  \foreach \x in {1.25,1.65,...,4.75}
    \draw[Brand,-{Stealth[length=1.7mm]}] (\x,0.62)--(\x,0.08);
  \draw[Brand] (1.2,0.62)--(4.8,0.62);
  \node[above] at (3,0.72) {\(q\)};
  \draw[<->] (0,-0.50)--node[below] {\(5a\)} (6,-0.50);
  \draw[<->] (1.2,-0.86)--node[below] {\(3a\)} (4.8,-0.86);
  \draw[<->] (0,-1.18)--node[below] {\(a\)} (1.2,-1.18);
  \node[below] at (0,0) {\(A\)};
  \node[below] at (1.2,0) {\(C\)};
  \node[below] at (4.8,0) {\(D\)};
  \node[below] at (6,0) {\(B\)};
  \figtitle{(i) 荷载关于跨中对称}
\end{tikzpicture}
\end{figure}
\end{problemcard}

\begin{solutioncard}
\textbf{答案：}\answer{\(F_A=F_B=1.5qa\) 向上。}

\textbf{解析：}均布荷载合力为
\[
  W=q\cdot3a=3qa,
\]
由于荷载段位于跨中，合力作用点距 \(A\) 为 \(2.5a\)。对 \(A\) 点取矩：
\[
  F_B\cdot5a-3qa\cdot2.5a=0,
\]
所以
\[
  F_B=\frac{3qa\cdot2.5a}{5a}=1.5qa.
\]
再由竖向平衡
\[
  F_A+F_B-3qa=0
  \quad\Rightarrow\quad
  F_A=1.5qa.
\]
也可直接利用对称性得到两端反力相等，各承担总荷载的一半。
\end{solutioncard}

\begin{problemcard}{2-1(j) 两个集中力偶作用的简支梁}
简支梁 \(AB\) 长 \(3a\)，\(C\) 点距 \(A\) 为 \(a\)。\(C\) 点作用力偶 \(M_e\)，\(B\) 点作用力偶 \(2M_e\)，方向如图。求两端竖向支反力。

\begin{figure}[H]
\centering
\begin{tikzpicture}[scale=0.9,>=Stealth]
  \draw[member] (0,0)--(5.4,0);
  \roller{0}{0}
  \roller{5.4}{0}
  \draw[Brand,-{Stealth[length=2mm]}] (1.75,0.70) arc[start angle=140,end angle=-40,radius=0.43];
  \draw[Brand,-{Stealth[length=2mm]}] (5.05,0.70) arc[start angle=140,end angle=320,radius=0.43];
  \node[above] at (1.8,0.78) {\(M_e\)};
  \node[above] at (5.0,0.78) {\(2M_e\)};
  \draw[<->] (0,-0.50)--node[below] {\(a\)} (1.8,-0.50);
  \draw[<->] (0,-0.86)--node[below] {\(3a\)} (5.4,-0.86);
  \node[below] at (0,0) {\(A\)};
  \node[below] at (1.8,0) {\(C\)};
  \node[below] at (5.4,0) {\(B\)};
  \figtitle{(j) 仅有力偶时两支反力方向相反}
\end{tikzpicture}
\end{figure}
\end{problemcard}

\begin{solutioncard}
\textbf{答案：}\answer{\(F_B=\dfrac{M_e}{3a}\) 向上，\(F_A=\dfrac{M_e}{3a}\) 向下。}

\textbf{解析：}全结构没有竖向外力，只有两个外力偶，因此两端竖向反力必然大小相等、方向相反。对 \(A\) 点取矩，按题解方向列式：
\[
  F_B\cdot3a+M_e-2M_e=0 .
\]
解得
\[
  F_B=\frac{M_e}{3a}.
\]
由竖向平衡
\[
  F_A+F_B=0
  \quad\Rightarrow\quad
  F_A=-\frac{M_e}{3a}.
\]
负号表示若假定 \(F_A\) 向上，则实际方向向下，大小为 \(\frac{M_e}{3a}\)。
\end{solutioncard}

\begin{problemcard}{2-1(k) 带中间铰的刚架}
图示刚架左柱 \(AD\) 承受水平均布荷载 \(q\)，总高度为 \(2a\)，上横梁中点 \(C\) 为铰。底部 \(A\) 与右侧 \(B\) 为支座，尺寸如图。求 \(A\)、\(B\) 处支座反力。

\begin{figure}[H]
\centering
\begin{tikzpicture}[scale=0.88,>=Stealth]
  \coordinate (A) at (0,0);
  \coordinate (D) at (0,3.2);
  \coordinate (C) at (2.0,3.2);
  \coordinate (E) at (4.0,3.2);
  \coordinate (B) at (4.0,1.6);
  \draw[member] (A)--(D)--(C)--(E)--(B);
  \node[pin] at (C) {};
  \roller{0}{0}
  \wallroller{4.0}{1.6}
  \foreach \y in {0.35,0.70,...,2.95}
    \draw[Brand,-{Stealth[length=1.7mm]}] (-0.55,\y)--(-0.05,\y);
  \draw[Brand] (-0.55,0.20)--(-0.55,3.10);
  \node[left] at (-0.60,1.65) {\(q\)};
  \draw[<->] (0,-0.58)--node[below] {\(a\)} (2,-0.58);
  \draw[<->] (2,-0.58)--node[below] {\(a\)} (4,-0.58);
  \draw[<->] (4.55,0)--node[right] {\(a\)} (4.55,1.6);
  \draw[<->] (4.55,1.6)--node[right] {\(a\)} (4.55,3.2);
  \node[below] at (A) {\(A\)};
  \node[above left] at (D) {\(D\)};
  \node[above] at (C) {\(C\)};
  \node[right] at (B) {\(B\)};
  \figtitle{(k) 取整体与 \(BC\) 部分隔离体}
\end{tikzpicture}
\end{figure}
\end{problemcard}

\begin{solutioncard}
\textbf{答案：}\answer{\(F_{By}=\dfrac23qa\)，\(F_{Bx}=\dfrac23qa\)，\(F_{Ay}=\dfrac23qa\)，\(F_{Ax}=\dfrac43qa\)。方向按图示反力方向标注。}

\textbf{解析：}左柱上的水平均布荷载合力为
\[
  q\cdot2a=2qa,
\]
作用在柱高的中点，距 \(A\) 的竖向距离为 \(a\)。取整体隔离体对 \(A\) 点取矩：
\[
  F_{By}\cdot2a+F_{Bx}\cdot a=q\cdot2a\cdot a .
\]
再取 \(BC\) 部分隔离体，对铰 \(C\) 点取矩，有
\[
  F_{By}\cdot a=F_{Bx}\cdot a,
\]
即
\[
  F_{By}=F_{Bx}.
\]
联立可得
\[
  3aF_{By}=2qa^2
  \quad\Rightarrow\quad
  F_{By}=F_{Bx}=\frac23qa .
\]
最后由整体平衡得到 \(A\) 处反力：
\[
  \sum F_y=0:\quad F_{Ay}=F_{By}=\frac23qa,
\]
\[
  \sum F_x=0:\quad F_{Ax}+F_{Bx}=2qa
  \quad\Rightarrow\quad
  F_{Ax}=2qa-\frac23qa=\frac43qa .
\]
\end{solutioncard}

\begin{problemcard}{2-1(l) 门式刚架局部均布荷载}
门式刚架宽 \(10\,\mathrm{m}\)，高 \(6\,\mathrm{m}\)，顶部中点 \(C\) 为铰。右半跨顶部受 \(2\,\mathrm{kN/m}\) 均布荷载，荷载长度 \(5\,\mathrm{m}\)。求两支座反力。

\begin{figure}[H]
\centering
\begin{tikzpicture}[scale=0.62,>=Stealth]
  \coordinate (A) at (0,0);
  \coordinate (D) at (0,6);
  \coordinate (C) at (5,6);
  \coordinate (E) at (10,6);
  \coordinate (B) at (10,0);
  \draw[member] (A)--(D)--(C)--(E)--(B);
  \node[pin] at (C) {};
  \wallroller{0}{0}
  \wallroller{10}{0}
  \foreach \x in {5.6,6.4,...,9.6}
    \draw[Brand,-{Stealth[length=1.8mm]}] (\x,7.05)--(\x,6.08);
  \draw[Brand] (5.0,7.05)--(10,7.05);
  \node[above] at (7.5,7.15) {\(2\,\mathrm{kN/m}\)};
  \draw[<->] (0,-1.0)--node[below] {\(5\,\mathrm{m}\)} (5,-1.0);
  \draw[<->] (5,-1.0)--node[below] {\(5\,\mathrm{m}\)} (10,-1.0);
  \draw[<->] (10.9,0)--node[right] {\(6\,\mathrm{m}\)} (10.9,6);
  \node[below] at (A) {\(A\)};
  \node[below] at (B) {\(B\)};
  \node[above] at (C) {\(C\)};
  \figtitle{(l) 右半跨顶部均布荷载}
\end{tikzpicture}
\end{figure}
\end{problemcard}

\begin{solutioncard}
\textbf{答案：}\answer{\(F_{By}=7.5\,\mathrm{kN}\)，\(F_{Bx}=\dfrac{25}{12}\,\mathrm{kN}\)，\(F_{Ay}=2.5\,\mathrm{kN}\)，\(F_{Ax}=\dfrac{25}{12}\,\mathrm{kN}\)。方向按题解图的支座反力方向。}

\textbf{解析：}右半跨均布荷载的合力为
\[
  W=2\times5=10\,\mathrm{kN},
\]
作用于右半跨中点，距 \(A\) 为 \(7.5\,\mathrm{m}\)。取整体对 \(A\) 点取矩：
\[
  F_{By}\times10=2\times5\times7.5,
\]
故
\[
  F_{By}=7.5\,\mathrm{kN}.
\]
再取 \(BC\) 部分隔离体对铰 \(C\) 点取矩。\(B\) 到 \(C\) 的水平距离为 \(5\,\mathrm{m}\)，竖向距离为 \(6\,\mathrm{m}\)，右半跨荷载合力到 \(C\) 的距离为 \(2.5\,\mathrm{m}\)，于是
\[
  F_{By}\times5=F_{Bx}\times6+2\times5\times2.5 .
\]
代入 \(F_{By}=7.5\,\mathrm{kN}\)，得
\[
  F_{Bx}=\frac{7.5\times5-25}{6}
  =\frac{12.5}{6}
  =\frac{25}{12}\,\mathrm{kN}.
\]
竖向平衡：
\[
  F_{Ay}+F_{By}=10
  \quad\Rightarrow\quad
  F_{Ay}=2.5\,\mathrm{kN}.
\]
水平平衡中没有外部水平荷载，因此两侧水平反力大小相等、方向相反，按题解方向均记大小为
\[
  F_{Ax}=F_{Bx}=\frac{25}{12}\,\mathrm{kN}.
\]
\end{solutioncard}

\begin{problemcard}{2-1(m) 矩形刚架侧向均布荷载}
矩形刚架宽 \(l\)、高 \(h\)。左柱 \(AB\) 受向右的水平均布荷载 \(q\)，底部 \(A\)、\(D\) 为支座。求 \(A\)、\(D\) 处支座反力。

\begin{figure}[H]
\centering
\begin{tikzpicture}[scale=0.9,>=Stealth]
  \coordinate (A) at (0,0);
  \coordinate (B) at (0,3.2);
  \coordinate (C) at (4.2,3.2);
  \coordinate (D) at (4.2,0);
  \draw[member] (A)--(B)--(C)--(D);
  \roller{0}{0}
  \wallroller{4.2}{0}
  \foreach \y in {0.45,0.85,...,3.05}
    \draw[Brand,-{Stealth[length=1.7mm]}] (-0.55,\y)--(-0.04,\y);
  \draw[Brand] (-0.55,0.25)--(-0.55,3.15);
  \node[left] at (-0.62,1.65) {\(q\)};
  \draw[<->] (0,-0.62)--node[below] {\(l\)} (4.2,-0.62);
  \draw[<->] (4.75,0)--node[right] {\(h\)} (4.75,3.2);
  \node[below] at (A) {\(A\)};
  \node[above left] at (B) {\(B\)};
  \node[above right] at (C) {\(C\)};
  \node[below] at (D) {\(D\)};
  \figtitle{(m) 侧向均布荷载刚架}
\end{tikzpicture}
\end{figure}
\end{problemcard}

\begin{solutioncard}
\textbf{答案：}\answer{\(F_{Dy}=\dfrac{qh^2}{2l}\) 向上，\(F_{Ay}=\dfrac{qh^2}{2l}\) 向下，\(F_{Dx}=qh\) 向左。}

\textbf{解析：}左柱水平均布荷载的合力为
\[
  W=qh,
\]
作用于柱高的中点，距 \(A\) 的竖向距离为 \(h/2\)。取整体对 \(A\) 点取矩：
\[
  F_{Dy}\cdot l=q h\cdot\frac{h}{2},
\]
故
\[
  F_{Dy}=\frac{qh^2}{2l}.
\]
竖向外荷载为零，所以两端竖向支反力大小相等、方向相反：
\[
  F_{Ay}+F_{Dy}=0
  \quad\Rightarrow\quad
  F_{Ay}=-\frac{qh^2}{2l}.
\]
因此 \(A\) 处实际竖向反力向下，大小为 \(\frac{qh^2}{2l}\)。水平方向平衡给出
\[
  F_{Dx}=qh,
\]
方向与水平分布荷载相反，即向左。
\end{solutioncard}
